% 来源及取材范围见分题文件；此为整理者转录的正文，不是下载原件。
\begin{theorem}[5.1: primitive element theorem]
Let $E=F[\alpha_1,\ldots,\alpha_r]$ be a finite extension of $F$, and assume that $\alpha_2,\ldots,\alpha_r$ are separable over $F$ (but not necessarily $\alpha_1$). Then there exists a $\gamma\in E$ such that $E=F[\gamma]$.
\end{theorem}
\begin{proof}
For finite fields, we proved this in 4.19. Hence we may assume $F$ to be infinite. It suffices to prove the statement for $r=2$, for then
\[F[\alpha_1,\alpha_2,\ldots,\alpha_r]=F[\alpha'_1,\alpha_3,\ldots,\alpha_r]=F[\alpha''_1,\alpha_4,\ldots,\alpha_r]=\cdots.\]
Thus let $E=F[\alpha,\beta]$ with $\beta$ separable over $F$. Let $f$ and $g$ be the minimal polynomials of $\alpha$ and $\beta$ over $F$, and let $L$ be a splitting field for $fg$ containing $E$. Let $\alpha_1=\alpha,\ldots,\alpha_s$ be the roots of $f$ in $L$, and let $\beta_1=\beta,\beta_2,\ldots,\beta_t$ be the roots of $g$. For $j\ne1$, $\beta_j\ne\beta$, and so the the equation
\[\alpha_i+X\beta_j=\alpha+X\beta\]
has exactly one solution, namely, $X=(\alpha_i-\alpha)/(\beta-\beta_j)$. If we choose a $c\in F$ different from any of these solutions (using that $F$ is infinite), then
\[\alpha_i+c\beta_j\ne\alpha+c\beta\quad\text{unless }i=1=j.\]
Let $\gamma=\alpha+c\beta$. I claim that $F[\alpha,\beta]=F[\gamma]$.

The polynomials $g(X)$ and $f(\gamma-cX)$ have coefficients in $F[\gamma]$, and have $\beta$ as a root:
\[g(\beta)=0,\qquad f(\gamma-c\beta)=f(\alpha)=0.\]
In fact, $\beta$ is their only common root, because we chose $c$ so that $\gamma-c\beta_j\ne\alpha_i$ unless $i=1=j$. Therefore
\[\gcd(g(X),f(\gamma-cX))=X-\beta.\]
Here we computed the gcd in $L[X]$, but this is equal to the gcd computed in $F[\gamma][X]$ (Proposition 2.17). Hence $\beta\in F[\gamma]$, and this implies that $\alpha=\gamma-c\beta$ also lies in $F[\gamma]$. This proves the claim.
\end{proof}

\paragraph{Separating transcendence bases.}
Let $E\supset F$ be fields with $E$ finitely generated over $F$. A subset $\{x_1,\ldots,x_d\}$ of $E$ is a \emph{separating transcendence basis} for $E/F$ if it is algebraically independent over $F$ and $E$ is a finite separable extension of $F(x_1,\ldots,x_d)$.

\begin{theorem}[9.27]
If $F$ is perfect, then every finitely generated extension $E$ of $F$ admits a separating transcendence basis over $F$.
\end{theorem}
\begin{proof}
If $F$ has characteristic zero, then every transcendence basis is separating, and so the statement becomes that of Theorem 9.10. Thus, we may assume $F$ has characteristic $p\ne0$. Because $F$ is perfect, every polynomial in $X_1^p,\ldots,X_n^p$ with coefficients in $F$ is a $p$th power in $F[X_1,\ldots,X_n]$:
\[\sum a_{i_1\cdots i_n}X_1^{i_1p}\cdots X_n^{i_np}
=\left(\sum a_{i_1\cdots i_n}^{1/p}X_1^{i_1}\cdots X_n^{i_n}\right)^p.\]
Let $E=F(x_1,\ldots,x_n)$, and assume that $n>d+1$, where $d$ is the transcendence degree of $E$ over $F$. After renumbering, we may suppose that $x_1,\ldots,x_d$ are algebraically independent (9.9). Then $f(x_1,\ldots,x_{d+1})=0$ for some nonzero irreducible polynomial $f(X_1,\ldots,X_{d+1})$ with coefficients in $F$. Not all $\partial f/\partial X_i$ are zero, for otherwise $f$ would be a polynomial in $X_1^p,\ldots,X_{d+1}^p$, which implies that it is a $p$th power. After renumbering $x_1,\ldots,x_{d+1}$, we may suppose that $\partial f/\partial X_{d+1}\ne0$.

Then $x_{d+1}$ is separably algebraic over $F(x_1,\ldots,x_d)$ and $F(x_1,\ldots,x_{d+1},x_{d+2})$ is algebraic over $F(x_1,\ldots,x_{d+1})$, hence over $F(x_1,\ldots,x_d)$ (1.31), and so, by the primitive element theorem (5.1), there is an element $y$ such that
\[F(x_1,\ldots,x_{d+2})=F(x_1,\ldots,x_d,y).\]
Thus $E$ is generated by $n-1$ elements (as a field containing $F$). After repeating the process, possibly several times, we will have $E=F(z_1,\ldots,z_{d+1})$ with $z_{d+1}$ separable over $F(z_1,\ldots,z_d)$.
\end{proof}
